On dit qu'un ensemble fini $\mathcal{S}$ de points du plan est \emph{équilibré} si, pour tous points $A$ et $B$
de $\mathcal{S}$ distincts, il existe un point $C$ de $\mathcal{S}$ tel que $AC = BC$. On dit que $\mathcal{S}$
est \emph{excentrique} si, pour tous points $A$, $B$ et $C$ de $\mathcal{S}$ distincts, il n'existe pas de
point $P$ de $\mathcal{S}$ tel que $PA = PB = PC$.
\begin{enumerate}
  \item[(a)] Prouver que pour tout entier $n \geqslant 3$, il existe un ensemble équilibré contenant exactement
  $n$ points.
  \item[(b)] Déterminer tous les entiers $n \geqslant 3$ pour lesquels il existe un ensemble équilibré et
  excentrique contenant exactement $n$ points.
\end{enumerate}
